\documentclass[11pt]{article}

% Manuscript v3 (draft, 17 Nov 2026). Every number below comes from a saved
% run; see Appendix A for the file behind each one. Items marked \pending{...}
% have NOT been run yet and must be filled or removed before any submission.
\usepackage[margin=1in]{geometry}
\usepackage{lmodern}
\usepackage{amsmath,amssymb}
\usepackage{booktabs}
\usepackage{graphicx}
\usepackage{array}
\usepackage{xcolor}
\usepackage{hyperref}
\hypersetup{colorlinks=true,linkcolor=blue,citecolor=blue,urlcolor=blue}
\emergencystretch=2em

\newcommand{\care}{\textsc{CARE}}
\newcommand{\evolve}{\textsc{ASEM-Evolve}}
\newcommand{\pending}[1]{\textcolor{red}{\textbf{[PENDING: #1]}}}

\title{When Does a Second Read Help?\\
Answer-Time Repair, Open Embedders, and Guarded Memory Evolution\\
for Long-Horizon Conversational QA}
\author{Anonymous authors}
\date{Draft v3}

\begin{document}
\maketitle

\begin{abstract}
Long-horizon conversational agents answer from histories that exceed their
context window. A common remedy is a bounded reader that may search again when
its first answer lacks support. We study such a reader, \care{}, which answers
from immutable, source-addressed turns and allows at most one repair read, and
we ask what limits it. On all 1,986 LoCoMo questions with a fixed Qwen3-8B
reader, repair raises the aggregate score by 0.0667 but reverses 46 correct
refusals; on all 500 LongMemEval cleaned-oracle questions it raises a diagnostic
answerable F1 from 0.348 to 0.374 while correct abstentions fall from 21/30 to
16/30. The same reader still refuses 16.7\% of answerable LoCoMo questions after
repair. Four further findings shape the design space. First, replacing the
MiniLM dense control with the open Qwen3-Embedding-0.6B reverses the earlier
lexical advantage at the retrieval level: annotated-evidence recall at eight
turns rises from 0.539 (BM25) to 0.617, and dense wins in all ten LoCoMo
conversations. Second, the refusals that repair destroys resist verification:
fourteen label-free trace features reach at most 0.60 AUC, and a Qwen3-8B yes/no
verifier lowers the net repair gain from +133.1 to +125.1 score points because
it never flags an unsupported premise. Third, most remaining LongMemEval errors
are reader-side aggregation failures: widening the evidence bundle raises
multi-session session recall from 0.74 to 0.93 on a 30-question probe without
changing F1. Fourth, the repair gate's risk term never triggers repair on its own,
and its speaker and event checks almost never fire outside LoCoMo. These results
motivate \evolve{}, a two-layer memory in which an immutable archive remains the
only citation source while a versioned derived layer is changed by guarded,
code-checked operators. We implement and unit-test the operators and
pre-register the evaluation; the evolution experiments themselves are not yet
run, and we make no claim that memory evolution improves accuracy.
\end{abstract}

\section{Introduction}
\label{sec:intro}

An agent that talks with a user over months must answer questions whose
evidence is scattered across sessions, may have been updated, and may be
attributed to the wrong person by a careless reader. Two families of systems
address this. \emph{Write-time} systems extract, link and consolidate memories
as the conversation arrives~\cite{packer2023memgpt,chhikara2025mem0,amem2025,memoryr12025,cao2026higmem}.
\emph{Read-time} systems keep the original history and decide, per question,
what to read~\cite{du2026memr3,wang2026cuemem,liu2026livectx,chen2026justmem,yu2026lazymem}.
A bounded corrective reader sits between them: it answers from a small bundle,
reports what evidence is missing, and is allowed one more read~\cite{jiang2023active,asai2023selfrag,yan2024crag}.

Our earlier manuscript introduced such a reader, \care{} (Conditional Answer
Repair with Source-Anchored Evidence), and showed that under an 8K budget it
outperforms a same-reader dense control on LoCoMo~\cite{maharana2024locomo}.
This version asks a different question: \emph{what limits a bounded repair
reader, and which of those limits can be removed by changing what is stored
rather than how it is read?} We answer with a set of controlled diagnostics on
two benchmarks and one fixed open-weights reader.

\paragraph{Findings.}
\begin{enumerate}
  \item \textbf{The repair trade-off replicates.} On LoCoMo and on all 500
  LongMemEval cleaned-oracle questions~\cite{wu2024longmemeval}, repair
  improves answerable questions and costs correct refusals
  (Section~\ref{sec:tradeoff}). The reader also refuses many answerable
  questions, a cost the earlier manuscript did not report.
  \item \textbf{A modern open embedder reverses the retrieval comparison.}
  With Qwen3-Embedding-0.6B~\cite{zhang2025qwen3embedding}, dense retrieval
  recovers more annotated evidence than BM25 on both benchmarks and in every
  LoCoMo conversation (Section~\ref{sec:retrieval}). The earlier
  ``sparse beats dense'' conclusion depended on a weak encoder.
  \item \textbf{Lost refusals are hard to verify.} Neither cheap trace features
  nor a same-model yes/no verifier separates correct repairs from repairs that
  answer an unsupported question (Section~\ref{sec:verify}).
  \item \textbf{The remaining error is mostly aggregation.} On multi-session
  questions, more recall does not yield more correct answers
  (Section~\ref{sec:errors}).
\end{enumerate}

\paragraph{Design.}
Findings 3 and 4 point away from further read-time tuning and toward storing
structure that a small reader cannot recompute reliably: versioned attribute
values, explicit links across sessions, and guarded updates. Section~\ref{sec:evolve}
describes \evolve{}, which keeps \care{}'s immutable archive as the only
citation source and evolves a separate derived layer through operators whose
preconditions are checked in code.

\paragraph{Status of each claim.}
Findings 1--4 are measured on saved runs. The \evolve{} operators are
implemented and unit-tested. The evolution protocol is pre-registered but not
yet executed; Section~\ref{sec:evolve} marks every missing result.

\section{Related Work}
\label{sec:related}

\paragraph{Write-time memory.}
MemGPT pages information between context and external storage~\cite{packer2023memgpt}.
Mem0 extracts and consolidates memories~\cite{chhikara2025mem0}. A-Mem builds
linked notes and revises their attributes~\cite{amem2025}; Memory-R1 learns
write operations with reinforcement learning~\cite{memoryr12025}. HiGMem keeps
a two-level event--turn memory whose event summaries are refreshed as new turns
arrive, and lets an LLM choose which linked turns to read~\cite{cao2026higmem}.
These systems change what is stored; none, to our knowledge, separates an
immutable citation archive from an evolving derived layer with code-checked
operator guards.

\paragraph{Read-time memory access.}
MemR$^3$ tracks evidence gaps while routing among retrieval, reflection and
answering~\cite{du2026memr3}. CueMem and LiveCtx rebuild query-specific
context from original turns~\cite{wang2026cuemem,liu2026livectx}. JustMem
extracts atomic memory cards once per session and selects, per query, a lookup,
a broadened composition, or a replay of source conversations; it does not update
cards after writing them~\cite{chen2026justmem}. LazyMem defers construction to
query time, labelling candidate messages keep or drop with a 4B model trained
by GRPO~\cite{yu2026lazymem}. \care{} is closest to this family; its distinctive
element is a one-repair ceiling with a reader that names missing evidence.

\paragraph{Corrective and adaptive retrieval.}
FLARE retrieves again when generation is uncertain~\cite{jiang2023active};
Self-RAG learns to retrieve and critique~\cite{asai2023selfrag}; Adaptive-RAG
chooses a retrieval strategy by question complexity~\cite{jeong2024adaptive};
CRAG grades retrieved evidence and triggers corrective actions~\cite{yan2024crag}.
Our verifier results (Section~\ref{sec:verify}) concern the step these methods
share: deciding whether a second read should override the first answer.

\paragraph{Evaluation.}
LoCoMo provides long two-speaker dialogues with category-labelled
questions~\cite{maharana2024locomo}; LongMemEval tests multi-session
reasoning, updates and abstention~\cite{wu2024longmemeval}. Published
LongMemEval accuracies use model judges and differ in setting; we report a
diagnostic token F1 that is not comparable to them. Claim-level grounding
checks such as RAGChecker~\cite{ru2024ragchecker} test a stronger notion of
support than our citation checks.

\section{The \care{} Read Path}
\label{sec:care}

A conversation is a sequence of immutable turns
$s_i=(\mathrm{id}_i,\mathrm{session}_i,t_i,p_i,x_i)$ with session date $t_i$,
speaker $p_i$ and text $x_i$. For a question $q$, the first pass ranks turns
with BM25 and serializes the top eight under a 10,000-character cap. A frozen
reader returns JSON with a short answer $a_1$, cited source IDs $C_1$ and at
most two searches $Q_1$ for missing evidence. Code then (i) drops citations
outside the visible bundle, (ii) turns an uncited answer into an abstention,
and (iii) applies speaker-binding and two event checks (holiday names and
cessation verbs). Repair is triggered by
\begin{equation}
\tau(q)=\mathbf{1}\bigl[r_1 \lor (a_1=\bot \land (Q_1\neq\varnothing \lor \rho(q)))\bigr],
\label{eq:gate}
\end{equation}
where $r_1$ marks an invalid first pass and $\rho$ flags temporal, update,
cross-speaker or plural-aggregation cues. The repair read unites unstemmed and
stemmed BM25 hits, session neighbours and hits for the proposed searches,
restores conversation order, applies the same cap and calls the reader once
more. Each question costs $1+\tau(q)$ model calls. Full details are in the
earlier manuscript; nothing in this section changed.

\section{Experimental Setup}
\label{sec:setup}

\paragraph{Benchmarks.}
\emph{LoCoMo:} ten conversations, 5,882 turns, 1,986 questions in five
categories (282, 321, 96, 841 and 446 questions). We follow the category names
used by A-Mem and HiGMem~\cite{amem2025,cao2026higmem}: 1 multi-hop, 2 temporal,
3 open-domain, 4 single-hop, 5 adversarial (unsupported premise). The earlier
manuscript labelled categories 1 and 4 as single-hop and conversational; 60\%
of category-1 references are comma-separated lists, consistent with multi-hop,
and the official scorer averages per-item F1 for that category. This mapping
should still be confirmed against the benchmark's own documentation.
\emph{LongMemEval cleaned oracle:} 500 questions over the annotated relevant
sessions only (948 sessions, 10,960 turns), including 30 abstention questions.
Because the oracle setting supplies the relevant sessions, it tests the reader
far more than retrieval.

\paragraph{Models.}
The reader is Qwen3-8B~\cite{yang2025qwen3} at a pinned revision, half precision
on two T4 GPUs through vLLM, 8,192-token window, temperature zero, non-thinking
mode, at most 224 generated tokens. The dense encoder is Qwen3-Embedding-0.6B
at a pinned revision, fp16, last-token pooling, instruction-prefixed queries
and documents of at most 1,024 tokens (no LoCoMo turn and one LongMemEval turn
were truncated). Every model is open-weights; no proprietary API is used.

\paragraph{Metrics.}
LoCoMo uses the benchmark scorer: per-item F1 for category 1, stemmed-token F1
for categories 2--4, and a refusal phrase match for category 5; the aggregate
is the mean per-question score. For LongMemEval we report a \emph{diagnostic}
normalized token F1 on answerable questions and phrase-based abstention on the
30 \texttt{\_abs} questions; this is not the official judged accuracy.
Annotated-evidence recall is the fraction of gold evidence IDs in the top-$k$
ranking or visible bundle; gold IDs never enter inference.

\paragraph{Statistics.}
LongMemEval questions have independent histories, so we bootstrap questions.
LoCoMo has ten conversations; we report per-conversation outcomes where they
matter, because percentile intervals from ten clusters are optimistic.

\section{Findings}
\label{sec:findings}

\subsection{Repair trades refusals for answers on both benchmarks}
\label{sec:tradeoff}

\begin{table}[t]
\centering
\caption{First pass versus conditional repair with the same reader. LoCoMo
values are the benchmark aggregate; LongMemEval values are diagnostic F1 on
answerable questions. Refusal rate counts answerable questions that receive an
abstention.}
\label{tab:tradeoff}
\small
\begin{tabular}{llrrrr}
\toprule
Benchmark & Policy & Score & Correct refusals & Refused answerable & Calls \\
\midrule
LoCoMo (1,986) & First pass & 0.4588 & 374/446 & 630/1,540 (40.9\%) & 1,986 \\
               & Repair (\care{}) & 0.5256 & 328/446 & 257/1,540 (16.7\%) & 2,990 \\
LongMemEval (500) & First pass & 0.348 & 21/30 & 195/470 (41.5\%) & 500 \\
                  & Repair & 0.374 & 16/30 & 156/470 (33.2\%) & 716 \\
\bottomrule
\end{tabular}
\end{table}

Table~\ref{tab:tradeoff} shows the same pattern on both benchmarks. On LoCoMo,
repair raises the aggregate by 0.0667 (conversation-cluster interval
[0.0503, 0.0848]); 279 answerable questions improve, none decline, and 46 of
446 correct refusals are lost. On LongMemEval, the answerable F1 rises by
0.026 (question bootstrap interval [0.014, 0.040]); 19 questions improve, none
decline, and 5 of 21 correct abstentions are lost with none gained.

Two quantities absent from the earlier manuscript matter as much as the
aggregate. First, the reader refuses many answerable questions: 16.7\% of
answerable LoCoMo questions remain refused after repair, and 33.2\% on
LongMemEval. Second, a policy that always refuses scores 446/1,986 = 0.2246 on
the LoCoMo aggregate, above the window-limited full-context baseline of the
earlier study (0.1961). Aggregate scores that mix answer F1 with refusal accuracy
should therefore be read with this floor in mind.

\paragraph{Repair is much less effective outside LoCoMo.}
On LoCoMo, 419 of 1,004 repairs (42\%) turn a refusal into an answer. On
LongMemEval, 39 of 195 repaired answerable questions (20\%) receive an answer,
and only 11 of those 39 reach F1 $\geq$ 0.5. The precision of repaired answers
(0.28) is far below that of first-pass answers that needed no repair (0.63,
$n=275$).

\paragraph{Repair can lower source recall.}
On LongMemEval, the share of annotated answer sessions in the visible bundle
falls after repair for multi-session (0.917 to 0.835) and temporal questions
(0.961 to 0.910). The expanded pool is restored to conversation order and cut
at the same character cap, so turns from later sessions can be dropped. We have
not isolated this mechanism.

\subsection{The gate's extra checks rarely act}
\label{sec:gate}

In the complete LoCoMo trace, every one of the 1,004 repairs is triggered by an
invalid first pass (137) or a non-empty search request (867); none is triggered
by the risk term $\rho$ alone. In Eq.~\eqref{eq:gate}, $\rho$ therefore never
changed a decision on this dataset. The speaker and event checks are written
for named two-speaker dialogue: among the 500 LongMemEval questions, none names
a holiday or a speaker label and one contains a cessation verb. Outside
LoCoMo, the operative check is the citation requirement.

\subsection{An open embedder reverses the retrieval comparison}
\label{sec:retrieval}

\begin{table}[t]
\centering
\caption{Recall of annotated evidence in the top-$k$ ranking (no reader
involved). LoCoMo: 1,531 answerable questions with evidence IDs; LongMemEval:
470 answerable questions, turn-level \texttt{has\_answer} labels. Hybrid is
reciprocal-rank fusion of unstemmed BM25 and dense ($k_{\mathrm{RRF}}=60$).}
\label{tab:retrieval}
\small
\begin{tabular}{lrrrrrr}
\toprule
 & \multicolumn{4}{c}{$k=8$} & \multicolumn{2}{c}{$k=18$} \\
\cmidrule(lr){2-5}\cmidrule(lr){6-7}
Benchmark & BM25 & BM25-stem & Dense & Hybrid & BM25 & Dense \\
\midrule
LoCoMo & 0.539 & 0.578 & \textbf{0.617} & 0.612 & 0.624 & \textbf{0.725} \\
LongMemEval & 0.837 & 0.846 & \textbf{0.926} & 0.911 & 0.941 & \textbf{0.988} \\
\bottomrule
\end{tabular}
\end{table}

The earlier manuscript found that BM25 beats dense MiniLM-L6-v2 retrieval at
eight turns (source recall 0.5353 versus 0.4576) and concluded that sparse
selection was the better first bundle. Table~\ref{tab:retrieval} replaces the
encoder with Qwen3-Embedding-0.6B. Dense recall at eight turns now exceeds
BM25 on LoCoMo (0.617 versus 0.539) and LongMemEval (0.926 versus 0.837), and
dense beats BM25 in all ten LoCoMo conversations (per-conversation gains
from +0.006 to +0.151). Fusion does not beat dense alone overall.

\begin{table}[t]
\centering
\caption{Recall at eight turns by question type.}
\label{tab:retrieval-type}
\small
\begin{tabular}{lrrrr}
\toprule
Type & $n$ & BM25 & Dense & Hybrid \\
\midrule
LoCoMo 1 multi-hop & 281 & 0.225 & \textbf{0.378} & 0.328 \\
LoCoMo 2 temporal & 320 & 0.648 & \textbf{0.754} & 0.720 \\
LoCoMo 3 open-domain & 89 & 0.238 & \textbf{0.368} & 0.329 \\
LoCoMo 4 single-hop & 841 & 0.634 & 0.671 & \textbf{0.695} \\
\midrule
LongMemEval knowledge-update & 72 & 0.912 & \textbf{0.975} & 0.961 \\
LongMemEval multi-session & 121 & 0.670 & \textbf{0.858} & 0.818 \\
LongMemEval temporal-reasoning & 127 & 0.827 & \textbf{0.910} & 0.890 \\
LongMemEval single-session-user & 64 & 0.953 & \textbf{1.000} & 0.984 \\
LongMemEval single-session-assistant & 56 & \textbf{0.982} & 0.929 & \textbf{0.982} \\
LongMemEval single-session-preference & 30 & 0.861 & \textbf{0.989} & 0.972 \\
\bottomrule
\end{tabular}
\end{table}

The largest gains fall on the question types that need several pieces of
evidence: LoCoMo multi-hop (0.225 to 0.378) and LongMemEval multi-session
(0.670 to 0.858), as Table~\ref{tab:retrieval-type} shows. These are recall
results only. We therefore ran the full matched-reader comparison with the
Qwen3-Embedding-0.6B top-8 as first bundle, keeping \care{}'s prompt, checks and
repair gate. Without repair, the dense first pass beats BM25 end to end (0.4730
versus 0.4588; paired $+0.0142$, conversation-cluster interval $[+0.0062,+0.0228]$,
9 of 10 conversations). With repair it scores 0.5169 against \care{}'s 0.5256
($-0.0087$, interval $[-0.0183,+0.0013]$, 3 of 10 conversations): answerable F1 is
equal (0.4640 versus 0.4648) and \care{} keeps 16 more correct refusals. The dense
system is better on multi-hop questions (0.3005 versus 0.2609) and uses fewer calls
(2,813 versus 2,990). Repair adds less on top of the dense bundle ($+0.0439$) than on
top of BM25 ($+0.0667$), consistent with a lexical repair expansion that complements a
lexical first bundle more.

\subsection{Lost refusals resist verification}
\label{sec:verify}

If a verifier could reject the repairs that answer unsupported questions while
keeping the useful ones, repair would lose its abstention cost. We study the
419 LoCoMo repairs that turn a first-pass refusal into an answer: 373 on
answerable questions (279 improve the score, 94 tie) and 46 on unsupported
questions, each of which loses a correct refusal. Accepting every flip adds
133.1 score points to the first-pass total, which is \care{}'s result. A perfect
verifier would add 179.1, a ceiling of +0.023 aggregate over \care{}.

\paragraph{Label-free trace features fail.}
We computed fourteen features available at inference, in two groups. Ten
describe the repair trace: number of searches, number of citations, whether a
citation is new relative to the first bundle and the fraction that are new,
number of new candidates, number of risk cues, answer length, numeric answer,
answer--question overlap and prompt length; their AUC for separating answerable
from unsupported flips lies between 0.48 and 0.60. Four measure lexical
coverage: coverage of the question's content words by the cited turns and by the
first bundle, the number of content words, and the absence of citations; their
AUC lies between 0.41 and 0.56. Coverage by the cited turns has AUC 0.46: wrong
flips cite turns that share \emph{more} words with the question. The reader appears
to cite a lexically related turn that does not establish the premise.

\paragraph{A same-model yes/no verifier fails.}
We asked Qwen3-8B, with the same pinned weights, whether the cited turns support
the proposed answer, offering three labels. It labelled 362 flips SUPPORTED, 57
NOT\_SUPPORTED and none PREMISE\_UNSUPPORTED. Accepting only SUPPORTED flips
keeps 325 of 373 answerable flips and catches 9 of 46 unsupported ones; the net
gain falls from +133.1 to +125.1, an aggregate of about 0.5215 instead of 0.5256.
The verifier is worse than no verifier because it never recognises an
unsupported premise, the exact failure it was meant to catch.

\paragraph{A quote-grounded verifier.}
We pre-registered a second verifier that must copy an exact premise quote and
an exact answer quote from the cited turns; code checks both quotes, and a flip
is accepted only if the verdict is SUPPORTED and both quotes are valid
(rule R3, thinking mode). Its success criterion is a net gain above +133.1.
It fails, with +90.6. In thinking mode the verifier labels 149 flips
PREMISE\_UNSUPPORTED and catches 35 of 46 unsupported flips, but it also rejects
185 of 373 useful ones. Across six scored configurations (two modes, four
acceptance rules) no rule beats accepting every flip; the best exploratory rule,
which rejects only PREMISE\_UNSUPPORTED verdicts, still loses 16.8 points.

\subsection{Where the remaining errors are}
\label{sec:errors}

\begin{table}[t]
\centering
\caption{LongMemEval diagnostic F1 by question type (first pass / repair) and
first-pass refusal rate on answerable questions. $n$ counts all questions of the
type, including abstention questions; F1 and refusal rate use answerable ones.}
\label{tab:lme-type}
\small
\begin{tabular}{lrrrr}
\toprule
Type & $n$ & First & Repair & Refused (first) \\
\midrule
single-session-user & 70 & 0.829 & 0.860 & 8\% \\
single-session-assistant & 56 & 0.676 & 0.755 & 12\% \\
knowledge-update & 78 & 0.541 & 0.555 & 15\% \\
temporal-reasoning & 133 & 0.190 & 0.220 & 51\% \\
multi-session & 133 & 0.075 & 0.082 & 69\% \\
single-session-preference & 30 & 0.009 & 0.014 & 77\% \\
\bottomrule
\end{tabular}
\end{table}

Table~\ref{tab:lme-type} locates the deficit in multi-session and temporal
questions. Preference questions have rubric-style references that token F1
cannot score, so their near-zero value says little about the system.

\paragraph{More evidence does not fix aggregation.}
A paired probe on 30 LongMemEval questions (12 multi-session, 9 temporal, 2
knowledge-update, 2 single-session-user, 5 abstention) compared three arms with
the same reader: the \care{} first bundle; a wider bundle (top 20 turns, 20,000
characters); and the wider bundle with a reader prompt that does not instruct
abstention on partial evidence. Widening raised multi-session session recall
from 0.74 to 0.93 and left answerable F1 unchanged (0.259 in both). The
relaxed prompt cut refusals from 60\% to 44\% but turned refusals into wrong
numbers (for example 350000 for a reference of \$25,000), left F1 at 0.250 and
lowered correct abstentions from 4/5 to 3/5. No multi-session question was
answered correctly in any arm. The probe is small and diagnostic, but it points
to counting, summing and ordering across turns as the bottleneck for an 8B
non-thinking reader.

\paragraph{Question dates help little.}
The LongMemEval runner had not passed \texttt{question\_date} to the reader. A
paired rerun of all 133 temporal questions reproduced the saved baseline exactly
(first pass 0.190, repair 0.220) and showed that adding the date raises repair
F1 by 0.0136 (10 questions better, 9 worse) and lowers refusals from 44\% to
40\%. The omission was a configuration error but does not explain the temporal
deficit.

\section{\evolve{}: Guarded Evolution of a Derived Memory Layer}
\label{sec:evolve}

The findings suggest that a bounded reader fails where answers must be
\emph{assembled}: counts across sessions, the current value of an attribute
that changed, or the order of two events. A small reader can see every piece and
still fail to combine them. Write-time structure can do that combination once,
in code, and expose it to the reader. The risk is that write-time memory
silently drifts or fabricates, which is why \care{} kept only raw turns.
\evolve{} is designed to obtain the first without the second.

\subsection{Two layers}

\emph{The archive} holds original turns and is append-only; it is the only
source a final answer may cite. \emph{The derived layer} holds attribute notes
$(\text{subject},\text{attribute},\text{value},\text{sources},\text{valid\_from},\text{version},\text{status})$,
typed links, and clusters. Every note cites archive turns, every change is
logged, and the derived layer can be rebuilt from the archive.

\subsection{Guarded operators}

\begin{table}[t]
\centering
\caption{Operators of the derived layer and the guards checked in code before
any change is applied. ADD, UPDATE and LINK are implemented and unit-tested;
the others are specified.}
\label{tab:ops}
\small
\begin{tabular}{p{1.6cm}p{5.2cm}p{6.6cm}}
\toprule
Operator & Purpose & Guard (rejected writes are logged, not applied) \\
\midrule
ADD & New attribute value & Cited turn exists; quote is an exact substring of it; value is a substring of the quote; no current note for the same key \\
UPDATE & Value changed over time & Same key has a current note; new evidence is not older than its \texttt{valid\_from}; value differs; old version kept as \emph{superseded} \\
LINK & Temporal, entity or causal edge & Kind in a closed set; no self-link; notes share an archive session or a cited turn \\
SPLIT / MERGE & Fix conflated or duplicated notes & Each part keeps its own source span; paraphrases are not counted as independent sources \\
PRUNE & Drop unused notes from the index & Never touches the archive; reversible \\
REORG & Re-cluster the layout & Previous layout stored for comparison and rollback \\
\bottomrule
\end{tabular}
\end{table}

Table~\ref{tab:ops} lists the operators. ADD, UPDATE and LINK are implemented
in the released code with fifteen unit tests covering archive immutability,
fabricated quotes, values absent from the quote, out-of-order updates,
unchanged values, version history, link guards, the audit log, and parsing of
both benchmarks' session-date formats (raw date strings are not chronologically
ordered in LoCoMo, so the update guard compares parsed dates).

\subsection{Signals for self-directed evolution}

Evolution is \emph{self}-directed only if no gold label drives it. The
candidate signals are those the read path already produces: a repair that finds
evidence after a refusal (a gap to fill), an invalid citation (a note to split
or update), two supported but conflicting values (a version to create), and
retrieval use of a note in accepted answers (a utility to raise by an
exponential moving average). Section~\ref{sec:verify} is a warning: on LoCoMo
the most natural signal, ``repair found an answer'', is right for answerable
questions and wrong for unsupported ones, and neither cheap features nor a
yes/no verifier can tell them apart. \evolve{} therefore must not write on a
repair signal until a verifier passes a pre-registered test on development data.

\subsection{Pre-registered evaluation}

The protocol streams sessions in date order and asks questions at fixed
checkpoints, on LongMemEval-S (full histories) as the test set and on LoCoMo
for development only. Arms: static archive (\care{}), append-only notes, full
evolution, evolution with each operator removed, and evolution driven by random
signals. Metrics: answerable F1 and correct abstentions per checkpoint;
stale-value, contradiction and duplicate rates in the derived layer; derived
layer size; write and read calls. Hypotheses: (H1) at equal read budget,
evolution beats the static archive on knowledge-update and multi-session
questions; (H2) label-free signals choose operators better than random ones;
(H3) contradiction and stale-value rates stay bounded as the stream grows. The
code and protocol will be frozen and tagged before the test run.
Results: \pending{H1--H3 on LongMemEval-S streams}.

\section{Discussion}
\label{sec:discussion}

\paragraph{What is supported.}
A bounded repair reader buys answer coverage with abstention accuracy on two
benchmarks, and leaves many answerable questions refused. A current open
embedder recovers more evidence than BM25, so the earlier preference for sparse
selection should not be generalised. The abstention cost of repair is not
removed by cheap features or a same-model yes/no check. On LongMemEval, the
remaining error is dominated by aggregation, not access.

\paragraph{What is not.}
We do not show that memory evolution improves accuracy; those experiments are
pending. The dense first bundle wins end to end only without repair; with repair
it is statistically indistinguishable from \care{} on this data. The
LongMemEval metric is a diagnostic token F1 in the oracle setting, not the
official judged accuracy on full histories. LoCoMo was used to develop \care{},
so its numbers are development results, and an independent audit of LoCoMo's
answer key has been reported that we have not yet incorporated. All results use
a single reader model.

\paragraph{Why this matters for memory design.}
Read-time systems such as JustMem and LazyMem avoid irreversible compression;
write-time systems such as HiGMem update summaries online~\cite{chen2026justmem,yu2026lazymem,cao2026higmem}.
Our diagnostics suggest a narrower role for writing: not compressing the
history, but storing the few derived facts that a small reader cannot combine
on its own, under guards that keep every stored fact traceable to an archive
turn.

\section{Conclusion}
\label{sec:conclusion}

We revisited a bounded repair reader for long-horizon conversational memory.
Repair helps answerable questions and hurts abstention on two benchmarks; an
open embedder overturns the earlier retrieval comparison; the refusals that
repair destroys cannot be recovered by the checks we tried; and the residual
error is aggregation. These findings motivate \evolve{}, a guarded two-layer
memory whose operators are implemented and whose evaluation is pre-registered.

\bibliographystyle{plain}
\bibliography{asem}

\appendix
\section{Provenance of Reported Numbers}
\label{app:provenance}

{\footnotesize
\begin{tabular}{>{\raggedright\arraybackslash}p{4.0cm}>{\raggedright\arraybackslash}p{11.5cm}}
\toprule
Result & Source \\
\midrule
LoCoMo \care{} traces, flips, over-refusal, gate triggers & \texttt{build/care/relation\_full/shard-*/\allowbreak .../report.json} (arm \texttt{bridge\_v2}) \\
LongMemEval 500 & \texttt{kaggle\_outputs/\allowbreak longmemeval\_20261116/\allowbreak shard-\{0..4\}} \\
Retrieval recall & \texttt{analysis/\allowbreak CARE\_embed\_retrieval\_qwen3\_0p6b\_20261117.json}; embeddings from kernel \texttt{asem-care-embed-qwen3-0-6b} \\
Label-free features & \texttt{analysis/\allowbreak ASEM\_Evolve\_flip\_signals\_20261116.txt}, \texttt{...coverage\_signal\_20261116.txt} \\
Yes/no verifier & \texttt{analysis/\allowbreak ASEM\_Evolve\_verifier\_result\_20261117.md} \\
30-question probe & \texttt{kaggle\_outputs/\allowbreak probe\_20261116} \\
Question-date rerun & \texttt{kaggle\_outputs/\allowbreak qdate\_20261117} \\
\evolve{} operators & \texttt{asem/evolve/}, \texttt{tests/test\_evolve\_p1.py} \\
\bottomrule
\end{tabular}}

\end{document}
